\documentclass[11pt]{article}

% --- arXiv-friendly preamble (standard packages only) ----------------------
\usepackage[utf8]{inputenc}
\usepackage[T1]{fontenc}
\usepackage{lmodern}
\usepackage[margin=1in]{geometry}
\usepackage{amsmath,amssymb}
\usepackage{booktabs}
\usepackage{microtype}
\usepackage{listings}
\usepackage{xcolor}
\usepackage{pgfplots}
\pgfplotsset{compat=1.17}
\usepackage[hidelinks]{hyperref}
\usepackage[french]{babel}

\lstset{basicstyle=\ttfamily\small,breaklines=true,frame=single,
        backgroundcolor=\color{gray!6},columns=fullflexible}

\title{\textbf{OctaSoma} : un octree 3-D efficace en cache pour la mémoire sémantique,\\
et les limites de la projection en très basse dimension}

\author{Projet OctaSoma\thanks{Auteur et affiliation à fixer avant
soumission. Sources, code et harnais de reproduction : le dépôt du projet.}}

\date{Juin 2026}

\begin{document}
\maketitle

\begin{abstract}
Les mémoires d'agent sont d'ordinaire opaques : un plongement est rangé dans un
index vectoriel en haute dimension et le rappel est une recherche du plus proche
voisin en boîte noire. Nous adoptons une posture différente. En partant de l'idée de
projeter les plongements en trois dimensions et de les indexer dans un
octree~\cite{octreerag}, nous traitons l'octree non comme un index plat mais comme
une \emph{mémoire fractale et multi-résolution} : parce que l'arbre est
auto-similaire à toute échelle, chaque profondeur est un \emph{niveau de zoom}, de
sorte qu'une requête peut être satisfaite -- et \emph{expliquée} -- à n'importe
quelle résolution, du thème large près de la racine au souvenir exact à une feuille.
\textbf{OctaSoma} est un moteur compact, à faible nombre de dépendances et écrit à
100\% en Rust sûr, qui réalise cela : une projection $3\times D$ apprise (PCA) ou
déterministe (Johnson--Lindenstrauss), un octree point-région à godets aux n\oe{}uds
de 64 octets, une recherche des $k$ plus proches voisins \emph{exacte} (prouvée
identique à la force brute, $38$--$77\times$ plus rapide qu'un balayage linéaire), un
rappel multi-résolution « zoomable », et -- parce que chaque souvenir porte de
vraies coordonnées 3-D -- une \emph{explicabilité et une visualisation natives}
qu'un index en haute dimension ne peut offrir. Nous mesurons aussi, honnêtement, le
coût de trois dimensions. Le rappel exact@1 s'effondre à ${\approx}0\%$, mais le
rappel \emph{thématique} (par cluster) est récupérable avec la PCA lorsque le corpus
comporte peu de thèmes dominants ($100\%$ à quatre thèmes, $73\%$ à seize, $3\%$ à
$256$ ; la projection aléatoire reste en retrait d'un facteur $4$--$6\times$). Une
comparaison directe avec un magasin exact en pleine dimension situe la niche
d'OctaSoma comme routeur \emph{grossier}, compact et explicable : des coordonnées
$85\times$ plus petites et, en pré-filtre, ${\sim}93\%$ du rappel par cluster en
pleine dimension à un coût de requête ${\sim}90\times$ moindre. Enfin, nous
présentons OctaSoma comme la fonction \emph{sémantique-spatiale} d'une pile de
mémoire d'agent en trois parties : une mémoire causale restreint le rappel à une
petite région, OctaSoma y fournit la couche sémantique bon marché et explicable, et
un noyau d'attention en consomme le résultat. Déployée \emph{par région} plutôt que
comme un index global unique, cette cascade récupère la cible sur des données réelles
de code d'agent avec un taux de succès de $99\%$ et ${\sim}26$ tokens par tour
(${\sim}137\times$ de moins qu'en injectant tout), là où l'index 3-D \emph{global}
d'OctaSoma seul obtient $0\%$ -- le routeur grossier n'est décisif qu'en tant que
brique de l'ensemble. Notre contribution est le modèle de mémoire multi-résolution
et explicable, son déploiement comme une fonction d'une pyramide d'inférence, ainsi
qu'une caractérisation honnête de sa plage de fonctionnement, livrés comme un
artefact propre et reproductible -- et non le pipeline de projection sous-jacent, que
nous attribuons aux travaux antérieurs~\cite{octreerag}.
\end{abstract}

\section{Introduction}
La récupération fondée sur les plongements sous-tend la mémoire sémantique moderne :
textes, images ou observations d'agents sont encodés comme des vecteurs dans
$\mathbb{R}^D$ (couramment $D\in[128,4096]$), et la pertinence est approchée par une
recherche des plus proches voisins. Les systèmes en production conservent le
plongement complet et s'appuient sur des index de plus proches voisins approchés
(ANN) tels que HNSW~\cite{hnsw} ou IVF-PQ~\cite{ivfpq,faiss} pour rendre cette
recherche praticable.

Cet article prend cet extrême opposé comme point de départ et en fait un \emph{modèle
de mémoire}. Au lieu de préserver les $D$ dimensions, nous réduisons le plongement à
\emph{trois} et l'indexons dans un octree~\cite{octreerag} ; mais plutôt que de
traiter l'arbre comme un index plat de plus proches voisins, nous exploitons sa
propriété caractéristique -- c'est un \emph{fractal}, auto-similaire à toute échelle
-- pour exposer toute la hiérarchie comme une \emph{mémoire zoomable et
multi-résolution}. Chaque profondeur est un niveau de zoom : un rappel peut être
parcouru du thème large près de la racine au souvenir exact à une feuille. Trois
dimensions rendent cela possible : un octree spatial devient naturel, les points sont
littéralement visualisables, les rappels deviennent \emph{explicables} (chaque
souvenir a de vraies coordonnées), et le stockage par élément est minimal. La
question est de savoir combien de sémantique utile un goulot aussi agressif conserve,
et sous quelles conditions.

Notre réponse est empirique et à deux faces. L'index lui-même est excellent : exact,
rapide, compact et sûr. La projection est le goulot d'étranglement, et nous la
quantifions précisément. Le véritable plus proche voisin d'un plongement en haute
dimension $D$ ne survit presque jamais à la projection en 3-D (rappel@1
$\approx 0\%$). Pourtant, la \emph{structure thématique grossière} survit, à condition
que (i) la projection soit \emph{apprise} (PCA, et non aléatoire) et que (ii) le
corpus comporte \emph{peu de thèmes dominants}. Nous donnons la courbe de
décroissance ainsi qu'une règle empirique.

\paragraph{Contributions.}
\begin{enumerate}
  \item Un \textbf{modèle de mémoire fractale et multi-résolution} : chaque niveau de
  l'octree sert de niveau de zoom, de sorte qu'un rappel peut être résumé et parcouru
  du thème grossier au souvenir exact (\texttt{zoom}/\texttt{zoom\_path}). À notre
  connaissance, cet usage navigable et multi-résolution d'un octree de plongements
  projetés est nouveau ; le travail le plus proche~\cite{octreerag} utilise l'arbre
  comme un index de récupération plat.
  \item Une \textbf{explicabilité et une visualisation natives} : parce que les
  souvenirs ont de vraies coordonnées 3-D, un rappel rend son « pourquoi » -- la
  position de la requête, les régions imbriquées qu'elle traverse et les souvenirs
  les plus proches avec leurs distances -- et le magasin s'exporte directement pour un
  visualiseur 3-D.
  \item Une \textbf{caractérisation honnête} de la projection en 3-D (rappel exact
  contre thématique, dimensionnalité intrinsèque, PCA contre aléatoire, une variante
  supervisée) et une \textbf{comparaison directe} avec un magasin en pleine dimension
  qui situe la niche de routeur compact / pré-filtre, avec un harnais entièrement
  reproductible.
  \item \textbf{OctaSoma}, une implémentation de référence écrite à 100\% en Rust sûr
  (\texttt{\#![forbid(unsafe\_code)]}) : recherche exacte des $k$ plus proches voisins
  en 3-D par séparation-évaluation, croissance dynamique du monde, persistance LZ4
  versionnée, un noyau d'agent et une CLI -- une seule dépendance d'exécution.
  \item Un \textbf{résultat système} : OctaSoma comme fonction sémantique-spatiale
  d'une pyramide d'inférence en trois parties (causal $\rightarrow$ sémantique
  $\rightarrow$ attention). Déployé \emph{par région causale} (un index « shardé »),
  le routeur 3-D autrement grossier aide à atteindre la cible avec $99\%$ de succès
  et ${\sim}26$ tokens par tour sur des données réelles -- une cascade où aucun
  composant seul ne suffit (Section~\ref{sec:pyramid}).
\end{enumerate}
Nous ne revendiquons explicitement \emph{pas} le pipeline plongements-vers-3-D-octree
comme nouveau ; il est dû à Ellendula et Bajaj~\cite{octreerag}.

\section{Travaux connexes}
\paragraph{Recherche vectorielle et ANN.} Les index à fichier inversé et à
quantification de produit~\cite{ivfpq,faiss} ainsi que les index de graphe tels que
HNSW~\cite{hnsw} dominent la récupération en haute dimension. Ils opèrent dans
l'espace d'origine (ou dans un espace légèrement compressé) et visent un rappel élevé
pour de grandes valeurs de $D$. OctaSoma n'est pas un concurrent ; il occupe l'extrême
opposé de l'axe de la dimensionnalité.

\paragraph{Réduction de dimensionnalité.} Le lemme de
Johnson--Lindenstrauss~\cite{jl} garantit que les applications linéaires aléatoires
préservent approximativement les distances par paires -- mais seulement vers des
dimensions logarithmiques en le nombre de points, bien au-delà de trois. L'analyse en
composantes principales~\cite{pca} capture les directions de variance maximale ; nous
en extrayons les trois premières par itération de la puissance avec déflation de
Hotelling~\cite{golubvanloan}. Les méthodes non linéaires (t-SNE~\cite{tsne},
UMAP~\cite{umap}) produisent de meilleures visualisations en 2-/3-D mais ne sont pas
de simples applications linéaires et ne fournissent pas d'index en ligne exact et peu
coûteux ; nous nous restreignons délibérément à une application linéaire $3\times D$.

\paragraph{Index spatiaux et conception orientée données.} Les octrees et les octrees
lâches~\cite{octree,ulrich} sont standards pour les requêtes de points et de régions
en 3-D. La recherche des plus proches voisins par séparation-évaluation avec bornes
admissibles est classique pour les arbres $k$-d et leurs apparentés~\cite{friedman}.
Notre implémentation suit la conception orientée données~\cite{dod} : des n\oe{}uds
POD contigus, référencés par index, dimensionnés à une ligne de cache.

\paragraph{Projection en basse dimension avec arbres spatiaux.} Le travail le plus
proche du nôtre, Ellendula et Bajaj~\cite{octreerag} regroupent les plongements de
phrases, projettent chaque groupe en trois dimensions et maintiennent des octrees
dynamiques par groupe pour une récupération hybride, rapportant des requêtes
environ $4{,}2\times$ plus rapides qu'un index à fichier inversé. OctaSoma partage
ce pipeline plongements-vers-3-D-octree ; nous ne le revendiquons pas comme
nouveau. Nous différons par la portée et l'accent : une unique projection
\emph{globale} PCA ou JL (nous mesurons PCA contre JL), un artefact exact,
reproductible et à dépendances minimales pensé pour la mémoire d'agent, et une
analyse explicite de \emph{quand} la projection échoue et survit. Nos résultats
rappel/dimension rejoignent la littérature ANN~\cite{drsurvey} : une réduction
agressive détruit le rappel du voisin exact tandis que l'appartenance à un groupe
reste significative~\cite{nnmeaning}, et les partitions spatiales alignées sur les
axes perdent leur avantage au-delà d'environ dix à vingt
dimensions~\cite{highdimann} --- ce qui explique précisément pourquoi un octree
impose une projection aussi agressive.

\section{Méthode}
\subsection{Projection vers $\mathbb{R}^3$}
Un plongement $x\in\mathbb{R}^D$ est appliqué par une matrice stockée par lignes
$W\in\mathbb{R}^{3\times D}$ vers $p = Wx \in \mathbb{R}^3$. Chaque coordonnée est un
produit scalaire accumulé en \texttt{f64} sur une boucle d'ordre fixe et de largeur 4,
ce qui rend $p$ \emph{identique au bit près d'une architecture à l'autre} -- essentiel,
car la disposition spatiale, et donc la réponse à chaque requête, doit être
reproductible.

Nous fournissons deux moyens d'obtenir $W$ :
\begin{itemize}
  \item \textbf{Johnson--Lindenstrauss (JL).} Remplir $W$ à partir d'un générateur
  Xorshift64 initialisé par une graine. Déterministe, instantané, indépendant des
  données.
  \item \textbf{PCA.} Centrer une matrice de calibration et extraire ses trois
  premières composantes principales par itération de la puissance avec déflation de
  Hotelling, en \texttt{f64}.
\end{itemize}
Dans les deux cas, les trois lignes sont normalisées en $L_2$. Pour un plongement de
norme unité, cela borne chaque coordonnée à $[-1,1]$ par Cauchy--Schwarz, de sorte
qu'un cube-monde unitaire contient déjà les données normalisées.

\subsection{Octree point-région à godets}
Les points sont indexés par un octree stocké comme un \texttt{Vec} contigu de n\oe{}uds
de 64 octets référencés par des indices \texttt{u32} (sans pointeurs, sans
\texttt{unsafe}) :
\begin{lstlisting}
struct OctreeNode {       // exactement 64 octets = 1 ligne de cache
  center:    [f32; 3],    // centre du cube
  half_size: f32,         // moitie de la longueur d'arete
  children:  [u32; 8],    // enfants par octant, NONE = absent
  bucket_id: u32,         // index du godet feuille, NONE = interne
  _padding:  [u8; 12],
}
\end{lstlisting}
Un n\oe{}ud est soit une \emph{feuille} (détenant une liste d'identifiants d'éléments
dans un tableau annexe \texttt{leaf\_buckets}), soit \emph{interne} (routant à travers
\texttt{children}). Garder les listes d'éléments hors du n\oe{}ud limite le parcours à
une ligne de cache par n\oe{}ud.

\paragraph{Insertion.} Pour insérer $p$, nous descendons depuis la racine en suivant
l'octant de $p$ (le bit $i$ est mis à 1 si et seulement si $p_i \ge c_i$). À une
feuille disposant de capacité libre, nous ajoutons ; une feuille pleine mais encore
divisible est \emph{subdivisée} (convertie en n\oe{}ud interne, son godet étant
redistribué entre huit enfants par octant) et la descente se poursuit. Comme chaque
subdivision pave exactement un cube en octants et que $p$ est garanti à l'intérieur du
cube racine, le confinement est préservé à chaque niveau -- l'invariant sur lequel
repose la recherche. Les points identiques au bit près partagent un godet feuille (un
vecteur extensible), de sorte que rien n'est jamais perdu.

\paragraph{Monde dynamique.} Si une coordonnée entrante dépasse la demi-arête
courante, le cube racine centré à l'origine \emph{double} jusqu'à ce qu'elle y tienne
et l'arbre est reconstruit à partir des points stockés. Le doublement rend la
croissance géométrique, donc amortie en $O(1)$ par insertion ; pour des plongements
normalisés, il ne se déclenche jamais.

\subsection{Recherche exacte des $k$ plus proches voisins}
Étant donné un point de requête $q$, \texttt{nearest}$(q,k)$ exécute une
séparation-évaluation : on visite un n\oe{}ud, et à un n\oe{}ud interne on descend
récursivement dans les enfants existants ordonnés par la borne inférieure
\begin{equation}
\underline{d}^2(q,\text{cube}) \;=\; \sum_{a\in\{x,y,z\}} \big(\max(0,\,|q_a-c_a|-h)\big)^2,
\end{equation}
la distance au carré de $q$ au point le plus proche d'un cube enfant (0 si à
l'intérieur), en \emph{élaguant} tout enfant dont la borne dépasse déjà la $k$-ième
meilleure distance courante. La borne étant admissible, aucun vrai voisin n'est jamais
écarté : \texttt{nearest} est \emph{exacte} -- identique au bit près à la force brute
-- et notre suite de tests l'affirme sur des milliers de points aléatoires. Les
recherches de charge utile top-1 et top-$k$ enveloppent cette primitive.

\subsection{Persistance et implémentation}
Le moteur sérialise vers un fichier \texttt{FRAC} v3 versionné et petit-boutiste dont
l'arène de charges utiles est compressée en LZ4 ; le chargeur valide le nombre magique,
la version et la dimensionnalité avant toute allocation, et ne panique jamais sur une
entrée malformée. L'ensemble de la \emph{crate} est du Rust sûr avec une seule
dépendance (un codec LZ4) ; les n\oe{}uds sont POD, le générateur aléatoire et la PCA
sont sans dépendance, et la projection est déterministe d'une plateforme à l'autre.

\section{Évaluation}
\label{sec:eval}
\subsection{Protocole}
Nous générons $N$ plongements de norme unité dans $\mathbb{R}^D$ à partir de $C$
clusters latents (« thèmes ») : chaque échantillon est un centre de cluster aléatoire
auquel s'ajoute un bruit gaussien renormalisé, ce qui confère une dimensionnalité
intrinsèque contrôlable via $C$. Les requêtes forment un ensemble \emph{tenu à
l'écart} issu des mêmes clusters avec un flux de bruit indépendant (interroger des
points insérés rendrait le rappel@1 trivialement égal à $100\%$). La vérité-terrain est
le plus proche voisin exact dans l'espace \emph{complet} en $D$ dimensions, par force
brute. Nous rapportons : le \emph{rappel exact@}$k$ (vrais voisins en haute dimension
$D$ récupérés par l'index 3-D), le \emph{rappel par cluster@}$k$ (éléments retournés
appartenant au thème de la requête -- la métrique dont la mémoire thématique a besoin),
la latence de requête (octree contre force brute sur les mêmes points 3-D ; les deux
exacts), le débit d'insertion et le ratio LZ4. Tous les chiffres proviennent d'une
seule machine, d'une compilation en mode release, et sont reproductibles via le
harnais fourni.

\subsection{Résultats}
Le tableau~\ref{tab:scale} rapporte le point de fonctionnement à grande échelle. La
recherche des $k$ plus proches voisins par octree est exacte et $67$--$77\times$ plus
rapide qu'un balayage 3-D linéaire ; la PCA atteint $70.8\%$ de rappel par cluster@1
là où une projection aléatoire ne parvient qu'à $13\%$ ; le rappel \emph{exact}@1 est
$\approx 0\%$ pour les deux. Les charges utiles se compressent d'un facteur
$5.5\times$.

\begin{table}[t]
\centering
\caption{Échelle : $N{=}50{,}000$, $D{=}256$, $C{=}16$ thèmes, $k{=}10$.
La latence et le débit dépendent de la machine.}
\label{tab:scale}
\begin{tabular}{lcccccc}
\toprule
Projection & cluster@1 & rappel@1 & $k$-NN octree & 3-D brute & accélération & insertions/s \\
\midrule
PCA (apprise) & \textbf{70.8\%} & 0.0\% & 5.69\,\textmu s & 383.6\,\textmu s & 67$\times$ & 1.86\,M \\
JL (aléatoire)   & 13.0\%          & 0.0\% & 5.12\,\textmu s & 393.6\,\textmu s & 77$\times$ & 1.53\,M \\
\bottomrule
\end{tabular}
\end{table}

La figure~\ref{fig:sweep} et le tableau~\ref{tab:sweep} montrent comment le rappel par
cluster dépend de la dimensionnalité intrinsèque ($N{=}20{,}000$, $D{=}128$). La PCA
préserve parfaitement la structure thématique pour un petit nombre de thèmes et
décline à mesure que les thèmes se multiplient ; trois axes principaux ne peuvent
séparer qu'une poignée de modes bien répartis. La JL reste loin derrière sur toute la
plage.

\begin{table}[t]
\centering
\caption{Rappel par cluster@1 en fonction du nombre de thèmes latents ($N{=}20{,}000$,
$D{=}128$, $k{=}10$). Le rappel exact@1 demeure $\approx 0\%$ partout.}
\label{tab:sweep}
\begin{tabular}{lcccc}
\toprule
thèmes $C$ & 4 & 16 & 64 & 256 \\
\midrule
PCA & \textbf{100.0\%} & \textbf{73.2\%} & 20.6\% & 2.6\% \\
JL  & 33.2\%           & 12.8\%          & 4.2\%  & 1.2\% \\
\bottomrule
\end{tabular}
\end{table}

\begin{figure}[t]
\centering
\begin{tikzpicture}
\begin{semilogxaxis}[
  width=0.82\linewidth, height=6cm,
  xlabel={nombre de thèmes latents $C$ (échelle log)},
  ylabel={rappel par cluster@1 (\%)},
  xtick={4,16,64,256}, xticklabels={4,16,64,256},
  ymin=0, ymax=105, log basis x=2,
  legend pos=north east, grid=major]
\addplot[mark=*,thick] coordinates {(4,100.0)(16,73.2)(64,20.6)(256,2.6)};
\addplot[mark=square*,thick,dashed] coordinates {(4,33.2)(16,12.8)(64,4.2)(256,1.2)};
\legend{PCA (apprise), JL (aléatoire)}
\end{semilogxaxis}
\end{tikzpicture}
\caption{Le rappel thématique se dégrade avec la dimensionnalité intrinsèque. Une
projection apprise (PCA) surpasse largement une projection aléatoire (JL), mais les
deux s'effondrent dès que le corpus comporte plus d'une poignée de thèmes dominants.}
\label{fig:sweep}
\end{figure}

\subsection{Analyse}
Quatre constats ressortent. (1) \textbf{L'index est exact et rapide} : l'élagage
procure un facteur $38$--$77\times$ sur un balayage linéaire et la marge croît avec
$N$. (2) \textbf{Trois dimensions détruisent les voisins exacts} : le rappel@1
$\approx 0\%$ est intrinsèque à la projection, non à l'index. (3) \textbf{La structure
thématique peut survivre}, mais seulement sous une projection apprise et une faible
dimensionnalité intrinsèque. (4) \textbf{PCA $\gg$ JL} pour cette tâche d'un facteur
$4$--$6\times$ ; si une passe de calibration est abordable, elle en vaut toujours la
peine.

\section{Discussion : quand trois dimensions suffisent}
OctaSoma se comprend au mieux comme un \emph{routeur sémantique grossier}. Avec une
projection PCA, il récupère de manière fiable les mémoires thématiquement pertinentes
lorsqu'un corpus comporte environ une douzaine de thèmes dominants ou moins -- un
régime qui couvre de nombreux cadres pratiques de mémoire d'agent (un petit ensemble
de personas, de tâches ou de domaines), l'indexation explicable ou visualisable (les
points sont littéralement en 3-D), le stockage compact embarquable, et le
pré-filtrage grossier en amont d'un récupérateur plus lourd. Il convient mal à la
recherche générale à fort rappel sur des corpus variés, où un index ANN en pleine
dimension est l'outil adéquat.

\section{OctaSoma comme une fonction d'une pyramide d'inférence}
\label{sec:pyramid}
Le verdict honnête de la Section~\ref{sec:eval} -- un index 3-D \emph{global} unique
est un routeur grossier dont le rappel exact s'effondre à mesure qu'un corpus se
diversifie -- n'est pas la fin de l'histoire mais une consigne de déploiement.
OctaSoma est conçu comme la couche \emph{sémantique-spatiale} d'une pile de mémoire
d'agent à trois fonctions, chacune détenant un type de mémoire différent :
\begin{itemize}
  \item \textbf{Causale / structurelle} (long terme) : un graphe de code/événements
  qui sait \emph{ce qui dépend de quoi}, mais rappelle de façon lexicale et non
  sémantique.
  \item \textbf{Sémantique / spatiale} (long terme) : OctaSoma -- rappel de
  plongements bon marché, explicable et visualisable.
  \item \textbf{Mémoire de travail compressée / attention} (court terme) : un noyau
  de cache KV quantifié qui score les tuiles mises en cache pour l'attention.
\end{itemize}
Les couches causale et d'attention sont des systèmes compagnons de la même
pile.\footnote{CCOS (causal) et SLHAv2 (attention) ; voir les dépôts du projet. La
contribution de cet article est OctaSoma et son rôle dans la cascade, non ces
systèmes.} CCOS indique lui-même que son rappel par tâche est « un point d'entrée
lexical délibérément simple\,\dots\ pas un récupérateur sémantique » -- précisément
la lacune que comble OctaSoma.

\paragraph{Déploiement par région.} La Section~\ref{sec:eval} montre que l'index 3-D
global échoue à l'échelle. Le remède est de le \emph{sharder} : conserver un petit
index OctaSoma par région causale (p.\,ex.\ un fichier source) au lieu d'un unique
index global. La couche causale restreint une requête à sa région (petit $N$) ;
OctaSoma classe alors \emph{à l'intérieur} de cette région, où trois dimensions
suffisent. Nous l'implémentons comme une mémoire shardée avec une projection PCA par
région (quelques lignes au-dessus du moteur de base).

\paragraph{La cascade, mesurée.} Le tableau~\ref{tab:cascade} exécute la pile sur
$795$ n\oe{}uds d'un code d'agent réel, avec des requêtes en langage naturel tenues à
l'écart et plongées par un modèle de dimension $768$. Lecture honnête :
\begin{itemize}
  \item \textbf{L'ensemble est décisif ; aucune brique seule ne suffit.} L'injection
  naïve est coûteuse ($3622$ tokens, $3\%$ pertinents) ; l'index 3-D global
  d'OctaSoma seul est à $0\%$ ; la région causale seule est correcte mais plus lourde.
  La cascade -- restriction causale puis rerank exact dans la petite région --
  atteint la cible $99\%$ du temps à ${\sim}26$ tokens par tour, ${\sim}137\times$ de
  moins que la naïve, avec un contexte entièrement pertinent sur le plan causal.
  \item \textbf{OctaSoma est la couche grossière, non le récupérateur précis.} Le
  succès est obtenu par la restriction causale plus le rerank exact ; le rôle
  défendable d'OctaSoma est la couche grossière bon marché, compacte, explicable et
  visualisable qui propose et organise -- en cohérence avec la Section~\ref{sec:eval},
  non en tension avec elle.
\end{itemize}

\begin{table}[t]
\centering
\caption{La cascade sur données réelles : $795$ n\oe{}uds d'un code d'agent en
production (sources de CCOS), $763$ requêtes tenues à l'écart, plongées par un modèle
de dimension $768$ (\texttt{nomic-embed-text}). « Pertinent (causal) » est la
fraction du contexte retourné située dans la région causale de la requête ; les
tokens sont par tour.}
\label{tab:cascade}
\begin{tabular}{lccc}
\toprule
stratégie & tokens/tour & cible atteinte & pertinent (causal) \\
\midrule
naïve (tout injecter)                      & 3622        & 100\%         & 3\%            \\
sémantique seule (OctaSoma 3-D global)     & 30          & \textbf{0\%}  & 4\%            \\
causale seule (région CCOS)                & 158         & 100\%         & 100\%          \\
\textbf{causale + sémantique (par région)} & \textbf{26} & \textbf{99\%} & \textbf{100\%} \\
\bottomrule
\end{tabular}
\end{table}

\paragraph{Visualiser le noyau d'attention.} OctaSoma sert aussi le noyau de mémoire
de travail comme une \emph{loupe}, non comme un routeur : projeter en 3-D les latents
de tuiles du cache KV (de dimension $128$) produit une carte navigable du cache,
colorée par tête, exposant les groupes de tuiles et la structure de compression. Nous
avons vérifié que ce pré-filtre 3-D ne récupère qu'une fraction du top-$k$
d'attention \emph{exact} ; la valeur honnête est donc l'inspection et la diversité --
non la sélection d'attention, dont le scoring propre du noyau a la charge.

\section{Un palier haute précision : empreintes SimHash}
\label{sec:sketch}
La projection 3-D est un routeur grossier par construction ; là où une précision au
voisin exact est nécessaire, nous ajoutons un palier bon marché et \emph{non
projectif}, plutôt qu'une projection non linéaire (qui sacrifierait l'index exact,
incrémental et explicable, sans pour autant préserver les angles fins en trois
dimensions). Une empreinte SimHash~\cite{simhash} du plongement \emph{complet} --
$b$ hyperplans aléatoires, un bit de signe chacun -- transforme la similarité cosinus
en distance de Hamming ($\mathbb{E}[\text{Hamming}]/b \approx \theta/\pi$), évaluée
par un unique \emph{population count}. Le classement par Hamming préserve bien plus
de structure angulaire que trois coordonnées ; utilisé comme \emph{présélection}
avant un reclassement cosinus exact sur les plongements stockés, il récupère la
plupart des vrais voisins que la projection écarte, à une fraction du coût d'un
balayage en force brute. Le tableau~\ref{tab:sketch} donne le rappel@$M$ (le voisin
exact parmi les $M$ candidats les moins chers ; le rappel@1 de l'hybride égale le
rappel@$M$ pour une présélection de taille $M$) sur les données de la
Section~\ref{sec:eval}.

\begin{table}[t]
\centering
\caption{Rappel@$M$ contre le plus proche voisin exact en pleine dimension
($N{=}20{,}000$, $D{=}768$, $16$ thèmes). La 3-D est indépendante des bits ; un
rappel@$M$ SimHash est le rappel@1 de l'hybride à présélection $M$.}
\label{tab:sketch}
\begin{tabular}{lcccc}
\toprule
méthode & @1 & @32 & @128 & @512 \\
\midrule
3-D PCA (octree) & 0.0\% & 3.0\% & 10.7\% & 46.7\% \\
SimHash-256 & 1.7\% & 12.0\% & 30.7\% & 70.3\% \\
SimHash-512 & 1.7\% & 17.0\% & 40.0\% & 82.3\% \\
SimHash-1024 & 1.7\% & 21.7\% & \textbf{52.7\%} & \textbf{88.7\%} \\
\bottomrule
\end{tabular}
\end{table}

Deux leviers échangent précision contre coût -- la largeur $b$ de l'empreinte et la
taille $M$ de la présélection (le reclassement est un produit scalaire par candidat).
Le palier est du Rust $100\%$ sûr et stable (un produit scalaire, un signe, et
\texttt{u64::count\_ones}) ; il échange de la mémoire contre de la précision et
\emph{complète} la couche 3-D, qui reste le routeur bon marché, explicable et
visualisable. C'est aussi le récupérateur \emph{global} précis de la pyramide
(Section~\ref{sec:pyramid}) lorsqu'une requête n'a pas de région causale.

\section{Limites}
La projection est linéaire par construction ; un apprentissage de variété non linéaire
pourrait relever le rappel thématique mais sacrifierait l'index en ligne exact et peu
coûteux. Nos clusters utilisent un bruit isotrope et sont plus accommodants que les
véritables variétés de plongements, de sorte que les chiffres de rappel par cluster
constituent une borne optimiste. La latence et le débit sont mono-machine et
mono-thread. Le moteur ne prend actuellement en charge que l'insertion (ni
suppression ni TTL), et la mutation prend \texttt{\&mut self}, laissant la politique
de concurrence à l'appelant.

\section{Conclusion}
Nous avons présenté OctaSoma, un petit octree 3-D exact en Rust sûr pour la mémoire
sémantique, et nous l'avons utilisé pour mesurer le point le plus extrême du compromis
dimensionnalité/fidélité. L'index spatial est exact, rapide et compact ; la projection
3-D est le facteur limitant, effaçant les voisins exacts mais -- sous PCA et faible
dimensionnalité intrinsèque -- conservant une structure thématique utile. La leçon
honnête est une règle empirique plutôt qu'une victoire au classement : ne réduisez à
trois dimensions que lorsque votre corpus est thématiquement étroit, apprenez la
projection, et traitez le résultat comme un routeur grossier. Son usage le plus fort
n'est donc pas autonome mais comme une fonction d'une pyramide d'inférence : shardé
par région causale, le routeur grossier devient décisif -- $99\%$ de rappel à
${\sim}26$ tokens par tour sur données réelles -- parce qu'une mémoire causale
restreint le problème et qu'un rerank exact le conclut. Le code et le harnais complet
de reproduction accompagnent cet article.

\begin{thebibliography}{99}
\bibitem{hnsw} Y.~Malkov and D.~Yashunin. Efficient and robust approximate nearest
neighbor search using Hierarchical Navigable Small World graphs. \emph{IEEE TPAMI},
2020.
\bibitem{ivfpq} H.~J\'egou, M.~Douze, and C.~Schmid. Product quantization for
nearest neighbor search. \emph{IEEE TPAMI}, 2011.
\bibitem{faiss} J.~Johnson, M.~Douze, and H.~J\'egou. Billion-scale similarity
search with GPUs. \emph{IEEE Trans.\ Big Data}, 2021.
\bibitem{jl} W.~Johnson and J.~Lindenstrauss. Extensions of Lipschitz mappings into
a Hilbert space. \emph{Contemp.\ Math.}, 1984.
\bibitem{pca} I.~T.~Jolliffe. \emph{Principal Component Analysis}. Springer, 2002.
\bibitem{golubvanloan} G.~Golub and C.~Van~Loan. \emph{Matrix Computations}. Johns
Hopkins University Press, 4th ed., 2013.
\bibitem{tsne} L.~van~der~Maaten and G.~Hinton. Visualizing data using t-SNE.
\emph{JMLR}, 2008.
\bibitem{umap} L.~McInnes, J.~Healy, and J.~Melville. UMAP: Uniform Manifold
Approximation and Projection. \emph{arXiv:1802.03426}, 2018.
\bibitem{octree} D.~Meagher. Geometric modeling using octree encoding.
\emph{Computer Graphics and Image Processing}, 1982.
\bibitem{ulrich} T.~Ulrich. Loose octrees. In \emph{Game Programming Gems}, 2000.
\bibitem{friedman} J.~Friedman, J.~Bentley, and R.~Finkel. An algorithm for finding
best matches in logarithmic expected time. \emph{ACM TOMS}, 1977.
\bibitem{dod} R.~Fabian. \emph{Data-Oriented Design}. 2018.
\bibitem{octreerag} A.~S.~Ellendula and C.~Bajaj. Self-Balancing, Memory-Efficient,
Dynamic Metric-Space Data Maintenance for Rapid Multi-Kernel Estimation. In
\emph{ECML-PKDD}, 2025. arXiv:2504.18003.
\bibitem{simhash} M.~S.~Charikar. Similarity estimation techniques from rounding
algorithms. In \emph{Proc.\ STOC}, 2002.
\bibitem{drsurvey} Dimensionality-reduction techniques for approximate nearest
neighbour search: a survey and evaluation. \emph{arXiv:2403.13491}, 2024.
\bibitem{nnmeaning} K.~Beyer, J.~Goldstein, R.~Ramakrishnan, and U.~Shaft. When is
``nearest neighbor'' meaningful? In \emph{ICDT}, 1999.
\bibitem{highdimann} W.~Li, Y.~Zhang, Y.~Sun, et al. Approximate nearest neighbor
search on high-dimensional data --- experiments, analyses, and improvement.
\emph{IEEE TKDE}, 2020.
\end{thebibliography}

\end{document}
